% ------------------------------------------------------------------------------
% Chapter 5
% ------------------------------------------------------------------------------
\chapter{Commands} 
\label{cha:commands}

A Makefile has been used to make the compilation and execution of the testbench.
Changing the parameter \texttt{TEST} it is possible to run 3 different tests as per request.

\section*{Compilation}

\texttt{comp} was used to make the compilation.
Important to note that, while the compilation flags are equal, the cover flag is used only for the dut.
This change remove the code coverage form the testbench.

\lstinputlisting[
    style=mkfile, 
    linerange={6-6}, 
    firstnumber=6
]{gcd/Makefile}

\lstinputlisting[
    style=mkfile, 
    linerange={12-16}, 
    firstnumber=12
]{gcd/Makefile}

\vspace{-2mm}
\section*{UVM Run}

\texttt{sim} was used to make the simulation and was configured using 4 parameters.
This snippet is also responsible for the creation of the waveforms as well as the coverage database.


\begin{itemize}
    \item \texttt{TEST} the name of the test to perform [ctrl\_brigup\_test, ctrl\_det\_test, ctrl\_rnd\_test]
    \item \texttt{VERBOSITY} the verbosity level of the log [UVM\_LOW, UVM\_MEDIUM, UVM\_HIGH]
    \item \texttt{SEED} the seed of the simulation [random, seed]
    \item \texttt{TIMEOUT} max simulation time allowed, units are required [example: 10000ns]
\end{itemize}

\lstinputlisting[
    style=mkfile, 
    linerange={0 - 4}, 
    firstnumber=0
]{gcd/Makefile}

\lstinputlisting[
    style=mkfile, 
    linerange={18 - 23}, 
    firstnumber=18
]{gcd/Makefile}

\section{Coverage Report}

\texttt{cov} was used to make the coverage report.

\lstinputlisting[
    style=mkfile, 
    linerange={26 - 28}, 
    firstnumber=26
]{gcd/Makefile}

\section*{FORMAL Run}

\texttt{formal} was used to execute the formal verification.

\lstinputlisting[
    style=mkfile, 
    linerange={30 - 31}, 
    firstnumber=30
]{gcd/Makefile}
